\documentclass[a4paper,11pt,parskip=half]{scrartcl}

\usepackage[T1]{fontenc}
\usepackage{lmodern}
\usepackage[ngerman]{babel}
\usepackage{csquotes}
\usepackage{amsmath}
\usepackage[top=22mm,bottom=25mm,left=25mm,right=25mm]{geometry}
\usepackage[dvipsnames,svgnames,x11names]{xcolor}
\usepackage{hyperref}
\usepackage[per-mode=fraction]{siunitx}
\usepackage{enumitem}
\usepackage{subcaption}
\sisetup{locale = DE}
\usepackage{schooltex}

\hypersetup{colorlinks, linkcolor=blue!50!black, urlcolor=blue!50!black}
\setlength{\parindent}{0pt}

% ------------------------------------------------------------
% Hilfsmittel nur für diese Dokumentation
% ------------------------------------------------------------
\tcbuselibrary{listings}
\lstset{
  basicstyle=\ttfamily\small,
  columns=fullflexible,
  keepspaces=true,
  breaklines=true,
  literate={ä}{{\"a}}1 {ö}{{\"o}}1 {ü}{{\"u}}1
           {Ä}{{\"A}}1 {Ö}{{\"O}}1 {Ü}{{\"U}}1 {ß}{{\ss}}1
           {’}{{'}}1,
}

% Ergebnis-Teil: \section wird zu \minisec, Tabellen/Raster passen in die Box
\tcbset{
  doku/.style={
    enhanced, breakable,
    colback=white, colframe=black!30, colbacklower=white,
    arc=1mm, boxrule=0.4pt,
    fonttitle=\bfseries\small, coltitle=black, colbacktitle=black!8,
    listing options={language={[LaTeX]TeX}, texcsstyle=*\color{blue!60!black}},
    before lower={\let\section\minisec\setlength{\textwidth}{\linewidth}},
    segmentation style={black!30, solid},
  },
}
\newtcblisting{beispiel}{doku, listing above text, title={Beispiel}}
\newtcblisting{beispielseite}{doku, listing side text, title={Beispiel}}
\newtcblisting{code}{doku, listing only, title={Code}}

\newcommand{\cs}[1]{\texttt{\textbackslash #1}}
\newcommand{\env}[1]{\texttt{#1}}
\newcommand{\marg}[1]{\texttt{\{}\meta{#1}\texttt{\}}}
\newcommand{\oarg}[1]{\texttt{[}\meta{#1}\texttt{]}}
\newcommand{\meta}[1]{\ensuremath{\langle}\textit{#1}\ensuremath{\rangle}}

\NewDocumentEnvironment{befehl}{m}
  {\begin{tcolorbox}[colback=black!6, colframe=black!6, boxrule=0pt, sharp corners,
     left=2mm, right=2mm, top=1mm, bottom=1mm, before skip=6pt, after skip=6pt,
     fontupper=\ttfamily\small]
   #1
   \end{tcolorbox}}
  {}

\newcommand{\hinweis}[1]{%
  \begin{tcolorbox}[colback=orange!6, colframe=orange!60!black, boxrule=0.4pt, arc=1mm,
    left=2mm, right=2mm, top=1mm, bottom=1mm]
    \small\textbf{Hinweis:} #1
  \end{tcolorbox}}

\title{Das Paket \texttt{schooltex}}
\subtitle{Dokumentation}
\author{Paul Baier}
\date{Stand: Oktober 2026}

\begin{document}
\maketitle
\tableofcontents

% ============================================================
\section{Überblick}
% ============================================================

\texttt{schooltex} bündelt Befehle und Umgebungen für den Schulalltag in zwei Bereichen:

\begin{itemize}
  \item \textbf{Unterrichtsskripte}: Gliederung, Tafelbilder, Versuchslisten mit
        Versuchsverzeichnis und Verweisen, Stundenverläufe und das Einbinden von Arbeitsblättern.
  \item \textbf{Arbeitsblätter und Klassenarbeiten}: Kopf, Aufgaben mit Punkten
        und automatischer Gesamtpunktzahl, Teilaufgaben, gesammelte Lösungen, Karoraster und
        Schlussblock.
\end{itemize}

Das Paket ist auf Deutsch ausgelegt. Die Ausgabetexte (\enquote{Aufgabe}, \enquote{Versuch},
\enquote{Viel Erfolg!}, \dots) sind fest eingebaut.

\subsection{Installation}

Das Repository wird direkt als persönlicher TeX-Baum geklont, also \emph{ohne} eigenen Projektordner:
\begin{code}
cd ~/texmf
git clone git@github.com:Ithu123/schooltex.git .
\end{code}
Danach findet \LaTeX{} das Paket ohne weitere Einstellungen. Die Icons in
\texttt{tex/latex/Icons} werden ebenfalls automatisch gefunden.

\subsection{Laden und Voraussetzungen}

\begin{code}
\documentclass{scrartcl}        % oder scrbook, scrreprt
\usepackage[ngerman]{babel}
\usepackage{csquotes}
\usepackage[bottom=2.5cm, top=22mm, right=2.5cm]{geometry}
\usepackage{hyperref}
\usepackage{siunitx}
\usepackage{enumitem}
\usepackage{subcaption}
\usepackage{xcolor}
\usepackage{schooltex}
\end{code}

\texttt{schooltex} lädt selbst unter anderem \texttt{tikz}/\texttt{pgfplots},
\texttt{tabularray}, \texttt{tcolorbox}, \texttt{pdfpages}, \texttt{graphicx}, \texttt{soul},
\texttt{caption}, \texttt{marginnote}, \texttt{xparse} und \texttt{cleveref}.
Einige Befehle setzen außerdem voraus, dass das Dokument Folgendes bereitstellt:

\begin{center}
\begin{tblr}{colspec={Q[l] X[l]}, hlines, row{1}={bg=black!8, font=\bfseries}}
  Voraussetzung & benötigt für \\
  KOMA-Klasse (\texttt{scrartcl}, \texttt{scrbook}, \dots) &
    Versuchsverzeichnis (\cs{DeclareNewTOC}), \cs{teilaufgabe} (\texttt{addmargin}) \\
  \texttt{enumitem} & Umgebung \env{experiments} \\
  \texttt{hyperref} & klickbare Verweise mit \cs{Vref} (ohne geht es auch, dann ohne Link) \\
  \texttt{siunitx} & \cs{header} (Einheit Minuten) \\
  \texttt{subcaption} & \cs{minicaption} \\
  \texttt{geometry} mit Rand & Punkte und Symbole am Rand bei \env{Aufgabe} \\
  \cs{Datum} & \cs{header} (muss im Dokument definiert werden) \\
  \cs{declarePath} & \cs{AB} \\
\end{tblr}
\end{center}

% ============================================================
\section{Unterrichtsskripte}
% ============================================================

\subsection{Gliederung: \cs{kap}, \cs{uber}, \cs{uuber}}

Drei eigene Gliederungsebenen mit unterstrichenen Überschriften, unabhängig von
\cs{section} \& Co. Die Sternform setzt die Überschrift ohne Nummer. Die Zähler hängen
voneinander ab: Ein neues \cs{kap} setzt \cs{uber} zurück usw.

\begin{befehl}{\textbackslash kap\{Titel\} \quad \textbackslash uber\{Titel\} \quad \textbackslash uuber\{Titel\} \quad \textit{jeweils auch mit} *}\end{befehl}

\begin{beispielseite}
\kap{Mechanik}
\uber{Kräfte}
\uuber{Gewichtskraft}
Text ohne Einzug.

\uber*{Ohne Nummer}
\end{beispielseite}

\subsection{Tafelbild: \cs{tafel}}

Eine grün beschriftete Box für den Tafelanschrieb. Sie darf über Seiten umbrechen
und Absätze enthalten.

\begin{befehl}{\textbackslash tafel\{Inhalt\}}\end{befehl}

\begin{beispiel}
\tafel{
  \textbf{Wellenlänge} $\lambda$: Abstand zweier benachbarter
  Punkte gleicher Phase.

  $c = \lambda \cdot f$
}
\end{beispiel}

\subsection{Versuche}

\subsubsection*{Umgebung \env{experiments} und \cs{V}}

Versuche werden in der Umgebung \env{experiments} gesammelt. Jeder \cs{V} beginnt einen
neuen, automatisch nummerierten Versuch (\enquote{V \meta{Abschnitt}.\meta{Nr}}).
Der Zähler beginnt in jedem \cs{section} neu.

\begin{befehl}{\textbackslash begin\{experiments\}[Überschrift] \dots\ \textbackslash end\{experiments\}\\
\textbackslash V \quad \textit{oder} \quad \textbackslash V[Name]}\end{befehl}

Mit einem \meta{Name} passiert dreierlei:
\begin{itemize}
  \item Der Name erscheint hinter der Nummer.
  \item Es wird das Label \texttt{V:\meta{Name}} gesetzt (für Verweise, siehe unten).
  \item Der Versuch landet im Versuchsverzeichnis.
\end{itemize}
Die Überschrift ist optional, Standard ist \enquote{Versuche}.

Zusätzlich gibt es zwei Hilfsbefehle für die Liste:
\begin{itemize}
  \item \cs{Vtext}\marg{Text} setzt einen Absatz ohne Nummer, etwa einen Hinweis zwischen zwei Versuchen.
  \item \cs{Vgroup}\marg{Titel} setzt eine blaue Zwischenüberschrift mit Linie, um Versuche zu gruppieren.
\end{itemize}

\begin{beispiel}
\begin{experiments}
  \Vgroup{Wellenwanne}
  \V[Wellenwanne eben] Ebene Wellen mit dem Linealerreger.
  \V[Wellenwanne Kreis] Kreiswellen mit dem Punkterreger.
  \Vtext{Beide Versuche am besten mit Stroboskop.}
  \V[Wellenwanne Spalt] Beugung am Spalt.
  \V Ein Versuch ohne Namen (kein Label, kein Eintrag im Verzeichnis).
\end{experiments}
\end{beispiel}

\subsubsection*{Verweise: \cs{Vref} und \cs{Vrefrange}}

Verweise nutzen den \emph{Namen} des Versuchs, das Präfix \texttt{V:} wird automatisch
ergänzt. Mehrere Versuche werden durch Kommas getrennt. Die Sternform schreibt
\enquote{Versuch} bzw.\ \enquote{Versuche} aus. Mit \texttt{hyperref} ist der ganze Verweis ein Link.

\begin{befehl}{\textbackslash Vref\{Name\} \quad \textbackslash Vref\{Name1, Name2, \dots\} \quad \textbackslash Vref*\{\dots\}\\
\textbackslash Vrefrange\{erster Name\}\{letzter Name\} \quad \textbackslash Vrefrange*\{\dots\}\{\dots\}}\end{befehl}

\begin{beispielseite}
\Vref{Wellenwanne eben}

\Vref*{Wellenwanne eben}

\Vref{Wellenwanne eben, Wellenwanne Spalt}

\Vref*{Wellenwanne eben, Wellenwanne Kreis, Wellenwanne Spalt}

\Vrefrange{Wellenwanne eben}{Wellenwanne Spalt}

\Vrefrange*{Wellenwanne eben}{Wellenwanne Spalt}
\end{beispielseite}

\hinweis{Weil das Komma Versuche trennt, darf ein Versuchsname selbst kein Komma enthalten.
Die alte Schreibweise \cs{Vref}\texttt{\{V:Name\}} funktioniert weiterhin.
Ein \cs{ref}\texttt{\{V:Name\}} gibt nur die Nummer aus, \cs{autoref}\texttt{\{V:Name\}} ergibt
\enquote{Versuch \meta{Nr}}.}

\subsubsection*{Versuchsverzeichnis}

Alle benannten Versuche werden gesammelt und mit \cs{listoftoc}\texttt{\{loe\}} ausgegeben,
typischerweise am Ende des Skripts:

\begin{beispiel}
\listoftoc*{loe}
\end{beispiel}
(\cs{listoftoc*} setzt die Liste ohne eigene Überschrift; ohne Stern erscheint \enquote{Versuchsverzeichnis}.)

\subsubsection*{Bilder zu Versuchen: \env{Vbilder} und \cs{Vbild}}

Mehrere Bilder nebeneinander, jedes in einer eigenen \texttt{minipage}. Das optionale Argument
von \env{Vbilder} ist die Standardbreite als Bruchteil der Zeilenbreite (Standard \texttt{0.45}).

\begin{befehl}{\textbackslash begin\{Vbilder\}[Breitenfaktor] \dots\ \textbackslash end\{Vbilder\}\\
\textbackslash Vbild[caption=\dots, label=\dots, width=\dots]\{Bilddatei\}}\end{befehl}

\begin{beispiel}
\begin{Vbilder}[0.3]
  \Vbild[caption={Aufbau}, label={fig:aufbau}]{example-image-a}
  \Vbild[caption={Messung}]{example-image-b}
  \Vbild[caption={Breiter}, width=0.35\linewidth]{example-image-c}
\end{Vbilder}
\end{beispiel}

\subsection{Teilbilder ohne Gleitumgebung: \cs{minicaption}}

Für Teilbilder in \texttt{minipage}s, ohne \env{figure}. Jedes Teilbild bekommt
\cs{minicaption}\oarg{Label}\marg{Text}, die Gruppe wird mit \cs{minimaincaption}\marg{Text}
abgeschlossen. Erst dieser Befehl setzt die Hauptbeschriftung und den Eintrag im Abbildungsverzeichnis.

\begin{beispiel}
\begin{center}
  \begin{minipage}{0.4\linewidth}\centering
    \includegraphics[width=\linewidth]{example-image-a}
    \minicaption[fig:links]{Linkes Teilbild}
  \end{minipage}\hfill
  \begin{minipage}{0.4\linewidth}\centering
    \includegraphics[width=\linewidth]{example-image-b}
    \minicaption{Rechtes Teilbild}
  \end{minipage}
  \minimaincaption{Zwei Teilbilder}
\end{center}
Siehe \ref{fig:links}.
\end{beispiel}

\subsection{Stundenverlauf: \env{plan}}

Eine fünfspaltige \texttt{tblr}-Tabelle mit den Spalten Zeit, Thema, Ausführung/Plan, Methode und
Medien/Material. Die Umgebung beginnt mit einem eigenen \cs{section}. Dessen Titel ist das optionale
Argument, Standard ist \enquote{Stundenverlauf}. Jede Zeile endet mit \verb|\\|.

\begin{befehl}{\textbackslash begin\{plan\}[Titel] Zeit \& Thema \& Ausführung \& Methode \& Material \textbackslash\textbackslash\ \dots\ \textbackslash end\{plan\}}\end{befehl}

Für die Methodenspalte gibt es Kürzel:

\begin{center}
\begin{tblr}{colspec={Q[l] Q[l] Q[l] Q[l] Q[l]}, hlines, row{1}={font=\ttfamily}}
  \textbackslash lsg & \textbackslash lv & \textbackslash ea & \textbackslash pa & \textbackslash ga \\
  \lsg & \lv & \ea & \pa & \ga \\
\end{tblr}
\end{center}

Mit \cs{umbruch}\marg{Zeile 1 \textbackslash\textbackslash\ Zeile 2} lassen sich Zeilenumbrüche
innerhalb einer Zelle erzwingen.

\begin{beispiel}
\begin{plan}
  10' & Einstieg & Schülerlaola, Frage nach dem Zusammenhang von
        $f$ und $\lambda$ & \lsg & \Vref{Wellenwanne eben} \\
  20' & Wellenfronten & Tafelaufschrieb anhand der Experimente
        & \umbruch{Demo \\ \pa} & \Vrefrange{Wellenwanne eben}{Wellenwanne Spalt} \\
\end{plan}
\end{beispiel}

\subsection{Arbeitsblätter einbinden: \cs{declarePath} und \cs{AB}}

\cs{AB} bindet das PDF eines Arbeitsblatts ins Skript ein und erzeugt einen Eintrag im
Inhaltsverzeichnis sowie das Label \texttt{pdf:\meta{Name}}. Erwartet wird die Ordnerstruktur
\texttt{\meta{Pfad}/\meta{Name}/\meta{Name}.pdf}, so wie sie beim Übersetzen eines Arbeitsblatts entsteht.
Den Pfad legt man einmal in der Präambel mit \cs{declarePath} fest.

\begin{befehl}{\textbackslash declarePath\{Pfad\} \\ \textbackslash AB[Modus]\{Name\}}\end{befehl}

\begin{center}
\begin{tblr}{colspec={Q[l] X[l]}, hlines, row{1}={bg=black!8, font=\bfseries}}
  Modus & Wirkung \\
  \texttt{default} (ohne Angabe) & jede Seite des PDFs als eigene Seite \\
  \texttt{small} & zwei Seiten verkleinert nebeneinander auf einer Querformatseite, mit Rahmen \\
  \texttt{smalllandscape} & zwei Querformatseiten verkleinert untereinander, mit Rahmen \\
  \texttt{landscape} & jede Seite im Querformat \\
\end{tblr}
\end{center}

\begin{code}
% Präambel
\declarePath{../../../Arbeitsblaetter/Physik/KlasseKS}
% im Text
\AB{KSUebung2DInferferenz}
\AB[small]{KSUebung2DInferferenz}
Siehe Arbeitsblatt auf Seite \pageref{pdf:KSUebung2DInferferenz}.
\end{code}

% ============================================================
\section{Arbeitsblätter und Klassenarbeiten}
% ============================================================

\subsection{Kopf: \cs{header}}

Graue Kopfzeile mit Bearbeitungszeit, Hilfsmitteln und Datum. Das Datum kommt aus dem Befehl
\cs{Datum}, den das Dokument selbst definieren muss.

\begin{befehl}{\textbackslash header\{Minuten\}\{Hilfsmittel\}}\end{befehl}

\begin{beispiel}
\newcommand{\Datum}{15.07.26}
\header{45}{Taschenrechner, Formelsammlung}
\end{beispiel}

\subsection{Aufgaben: \env{Aufgabe}}

\begin{befehl}{\textbackslash begin\{Aufgabe\}[Punkte] \dots\ \textbackslash end\{Aufgabe\}}\end{befehl}

Setzt die Überschrift \enquote{Aufgabe \meta{Nr}} mit fortlaufender Nummer.
Im Seitenrand erscheint links ein Symbol für \enquote{schriftlich}.
Werden \meta{Punkte} angegeben, steht rechts außerdem \enquote{\rule{0.8cm}{0.5pt}/\meta{Punkte}}.
Dezimalpunkte wie \texttt{3.5} sind erlaubt. Alle Punkte werden automatisch zur Gesamtpunktzahl
\cs{Punkte} aufsummiert.

\begin{code}
\begin{Aufgabe}[4]
  Benenne die drei Newtonschen Axiome und erkläre eines davon.
\end{Aufgabe}

\begin{Aufgabe}[3.5]
  Berechne die Gewichtskraft eines Körpers mit $m = \SI{2}{kg}$.
\end{Aufgabe}
\end{code}

\hinweis{Die Umgebung hat noch ein zweites optionales Argument für andere Randsymbole
(Taschenrechner, ohne Taschenrechner, Sterne für Niveaustufen). Diese Auswahl funktioniert
derzeit nicht: Jeder Wert außer \texttt{0} führt zu einem Fehler. Ältere Dokumente definieren dafür noch
\cs{schriftlich}, \cs{mitTR}, \cs{ohneTR}, \cs{einS}, \cs{zweiS} und \cs{dreiS}.}

\subsection{Teilaufgaben: \cs{teilaufgabe}}

Nummeriert mit a), b), c), \dots\ und leicht eingerückt. Die Zählung beginnt in jeder Aufgabe
und in jeder Lösung neu.

\begin{befehl}{\textbackslash teilaufgabe\{Text\}}\end{befehl}

\begin{beispiel}
\teilaufgabe{Berechne den Gangunterschied $\delta$.}
\teilaufgabe{Entscheide, ob konstruktive oder destruktive Interferenz vorliegt.}
\end{beispiel}

\subsection{Lösungen: \cs{solution} und \cs{printsols}}

\cs{solution} steht \emph{innerhalb} einer Aufgabe und speichert ihre Lösung, ohne sie auszugeben.
Die Lösung darf Absätze und Teilaufgaben enthalten. Mit \cs{printsols} werden die gespeicherten
Lösungen später ausgegeben, meist am Dokumentende oder auf einer eigenen Seite.

\begin{befehl}{\textbackslash solution\{Lösung\} \\
\textbackslash printsols \quad \textbackslash printsols[Nummer] \quad \textbackslash printsols[last]}\end{befehl}

\begin{itemize}
  \item \cs{printsols} gibt alle Lösungen aus.
  \item \cs{printsols}\texttt{[2]} gibt nur die Lösung zu Aufgabe 2 aus.
  \item \cs{printsols}\texttt{[last]} gibt nur die Lösung zur zuletzt gesetzten Aufgabe aus,
        praktisch direkt nach der Aufgabe.
\end{itemize}

\begin{code}
\begin{Aufgabe}[2]
  \teilaufgabe{Nenne die Formel für die Wellengeschwindigkeit.}
  \teilaufgabe{Berechne $c$ für $\lambda = \SI{2}{m}$ und $f = \SI{5}{Hz}$.}
  \solution{
    \teilaufgabe{$c = \lambda \cdot f$}
    \teilaufgabe{$c = \SI{2}{m} \cdot \SI{5}{Hz} = \SI{10}{m/s}$}
  }
\end{Aufgabe}
...
\newpage
\printsols
\end{code}

\subsection{Karoraster: \cs{grid}}

Ein \SI{0,5}{cm}-Karoraster über die volle Breite in ganzen Zentimetern. Ohne Argument füllt
das Raster den Rest der Seite, mit Argument hat es die angegebene Höhe in Zentimetern.

\begin{befehl}{\textbackslash grid \quad \textbackslash grid[Höhe in cm]}\end{befehl}

\begin{beispiel}
\grid[2]
\end{beispiel}

\subsection{Schluss: \cs{Ende}}

Setzt \enquote{Viel Erfolg!} und eine Tabelle für Gesamtpunktzahl, Note und Note in Worten.
Die Gesamtpunktzahl \cs{Punkte} wird automatisch aus allen Aufgaben summiert. Da sie beim
Durchlauf berechnet wird, muss \cs{Ende} \emph{nach} den Aufgaben stehen.

\begin{beispiel}
\Ende
\end{beispiel}
(Im Beispiel steht \enquote{0}, weil diese Dokumentation keine Aufgaben mit Punkten setzt.)

\subsection{Gerüst einer Klassenarbeit}

\begin{code}
\documentclass[a4paper]{scrartcl}
\input{Prambel/Prambel}
\newcommand{\Datum}{15.07.26}

\begin{document}
\header{60}{Taschenrechner}

\begin{Aufgabe}[4]
  Benenne die drei Newtonschen Axiome und erkläre eines davon.
  \solution{...}
\end{Aufgabe}
\grid

\begin{Aufgabe}[6]
  \teilaufgabe{...}
  \teilaufgabe{...}
  \solution{\teilaufgabe{...} \teilaufgabe{...}}
\end{Aufgabe}
\grid[8]

\Ende

\newpage
\printsols
\end{document}
\end{code}

% ============================================================
\section{Weitere Befehle}
% ============================================================

\begin{center}
\begin{tblr}{colspec={Q[l,font=\ttfamily] X[l]}, hlines, row{1}={bg=black!8, font=\bfseries}}
  {\normalfont Befehl} & Bedeutung \\
  \textbackslash Punkte & aktuelle Gesamtpunktzahl aller bisherigen Aufgaben \\
  \textbackslash Path & der mit \cs{declarePath} gesetzte Pfad \\
  \textbackslash Vbildbreite & Länge: Standardbreite von \cs{Vbild} (wird von \env{Vbilder} gesetzt) \\
  \textbackslash getpagetotal, \textbackslash gettextheight & interne Hilfsbefehle für \cs{grid} \\
  Zähler kap, unterkap, unterunterkap & Zähler für \cs{kap}, \cs{uber}, \cs{uuber} \\
  Zähler experiment & Versuchszähler, \cs{theexperiment} $=$ \meta{Abschnitt}.\meta{Nr} \\
  Zähler Aufgabe, TA & Aufgaben- und Teilaufgabenzähler \\
\end{tblr}
\end{center}

\end{document}
